\section{Problem Statement}
\label{sec:problem}

\subsection{The Communication Barrier}

People who communicate in Pakistan Sign Language (PSL) and people who do not know it often have no shared language. Published work on PSL describes this as a communication gap between the deaf and non-deaf communities~\cite{imran2021psl}. In everyday settings such as a classroom, a shop, a clinic or a family gathering, a conversation between a signer and a non-signer then depends on a human interpreter, on writing things down, or on improvised gestures. Interpreters are not always available. Writing is slow and depends on both people sharing a written language. Gestures break down as soon as the message becomes specific, for example a person's name or an address.

Non-signers who want to meet signers halfway, such as relatives, classmates and teachers, need a way to learn and practise signs, and to check whether they are reading or producing a sign correctly.

\subsection{Limited PSL Data}

Publicly available PSL datasets are limited~\cite{colorglove_psl}. The ones we found differ in what they contain: one captured 40 images for each hand configuration of the Urdu alphabet~\cite{imran2021psl}; UAlpha40 contains 36 static signs as JPG images and 4 dynamic signs as MP4 videos~\cite{sajid2025ualpha40}; and a Kaggle dataset contains 37 Urdu alphabet signs performed by 9 non-expert subjects~\cite{butt_psl_kaggle}. The large, labelled alphabet dataset we could download for a baseline is an American Sign Language (ASL) dataset~\cite{asl_alphabet_kaggle}, not a PSL one. A model trained only on ASL cannot be used for PSL, because the two are different sign languages with different signs.

\subsection{Why Start with Letters}

In PSL, fingerspelling represents the letters of the Urdu alphabet~\cite{colorglove_psl}, and it is used for proper names, technical terms, addresses and uncommon words~\cite{pslcfrt}. A tool that recognises individual letters therefore helps exactly in the situations where improvised gestures fail.

Letters are also a sensible first step for a technical reason. Each static letter is a single hand shape that can be captured in one image, the set of classes is small and fixed, and every prediction can be checked against a known answer. Recognising words, sentences or signs that involve movement needs sequences of frames and a much larger dataset, which is outside what one course project can build and evaluate properly.

\subsection{Which Letters}

There is no single agreed list of PSL letters to recognise. The published sources above report different numbers of signs: 36 static signs~\cite{sajid2025ualpha40}, 34 static alphabets~\cite{colorglove_psl} and 37 alphabet signs~\cite{butt_psl_kaggle}. We have not yet verified a letter list ourselves. The PSL letter set will therefore be fixed in Deliverable~2 from a cited source and checked with a PSL user. The current label map, with the 26 letters A to Z, fits the ASL baseline only.

\subsection{Scope}

This project covers \textbf{static letter signs only}. The system receives a single image of one hand and returns the most likely letters with their confidence. It does not recognise motion-based signs, words or sentences, and it does not translate speech or text into signs.
